% homework/hw03x/hw03x.tex
% Year 7 - Extra Homework: multiplication and division.
%
% Goes out alongside Homework 3, for students who will get more out of Unit 1
% with more practice on the two operations the whole unit runs on. The framing
% matters and is deliberate: this is not presented as catch-up work, because it
% is not. Factors, the HCF and the divisibility tests are all built on
% multiplying and dividing, so the cover says exactly that -- these are the two
% skills that make the rest faster, and this is extra practice on them.
%
% Nothing in here is labelled easy, basic or extra help, and no question is
% marked as being for anybody in particular. The sequence goes tables ->
% multiplication -> division -> division with remainders -> missing numbers ->
% "use it" (which lands straight back on Homework 3's factor questions) ->
% word problems. E6 deliberately re-uses two divisions the student has already
% done in E3 and E4, so the last section rewards looking back.
%
% House style is shared: see homework/hw-style.tex. Method: the new-homework
% skill. Source is pure ASCII on purpose.
%
% Build:  pwsh .claude/skills/new-homework/scripts/build-hw.ps1 hw03x

\documentclass[12pt,a4paper]{article}

% -- Answer key switch -------------------------------------------------------
\newif\ifkey
\ifdefined\TEACHER\keytrue\else
  \keyfalse        % \keyfalse = student copy   |   \keytrue = teacher copy
\fi

% -- House style -------------------------------------------------------------
% homework/hw-style.tex
% Shared house style for every Year 7 homework packet. \input this from the
% packet file, right after \documentclass and the answer-key switch.
%
% Do not fork this file per packet. If a packet needs a different look, the
% question is whether every packet should have it.
%
% The choices here are deliberate and were arrived at by printing and looking:
%   * No solid dark banners. Every header is a light tint with a coloured spine
%     on the left, so colour reads clearly without flooding a school printer.
%   * No marks and no mark totals. These are practice, not exam papers.
%   * Lato at 12pt with open leading, plus mathastext so inline maths is Lato
%     too. Without mathastext every "-6 + 4" is Computer Modern serif and the
%     sheet looks half-typeset.
%   * Writing lines are spaced for Year 7 handwriting, not adult margin notes.
%
% Provides: \wlines \blank \tickbox \sep \qmk-free marks (none) \hwsection
% \qhead \expr, the workspace box, and the remember / wordhelp / activity /
% yourturn coloured boxes. Palette matches the lesson decks exactly.

% -- Packages ----------------------------------------------------------------
\usepackage[T1]{fontenc}
\usepackage[utf8]{inputenc}
\usepackage[default]{lato}
\usepackage[a4paper,top=1.7cm,bottom=1.7cm,left=2.0cm,right=2.0cm,headheight=15pt]{geometry}
\usepackage{amsmath,amssymb}
\usepackage{xcolor}
\usepackage[most]{tcolorbox}
\usepackage{enumitem}
\usepackage{tikz}
\usetikzlibrary{arrows.meta}
\usepackage{array}
\usepackage{tabularx}
\usepackage{fancyhdr}
\usepackage{multicol}
\usepackage{needspace}
\usepackage{graphicx}

% Make maths use Lato too. Without this, every "-6 + 4" on the page is set in
% Computer Modern serif and the sheet looks half-typeset. Must come last.
\usepackage[italic]{mathastext}

\linespread{1.06}\selectfont

% -- The Learner's Book palette (same hex values as the lesson decks) ---------
\definecolor{cbteal}{HTML}{0087A8}
\definecolor{cbpurple}{HTML}{5C2483}
\definecolor{cborange}{HTML}{C25E12}
\definecolor{cbgreen}{HTML}{4A8B23}
\definecolor{cbred}{HTML}{C8102E}
\definecolor{cbblue}{HTML}{1A5FA8}
\definecolor{rulegray}{HTML}{9AA5AD}

% -- Answer lines and blanks -------------------------------------------------
% \wlines{n} draws n ruled writing lines, spaced wide enough for Year 7
% handwriting rather than for an adult with a fine pen.
\makeatletter
\newcommand{\wlines}[1]{%
  \par\vspace{0.75em}%
  \@tempcnta=0\relax
  \@whilenum\@tempcnta<#1\do{%
    \hbox to \linewidth{\color{rulegray}\leaders\hrule height 0.5pt\hfill}%
    \vspace{1.5em}%
    \advance\@tempcnta by 1\relax}%
  \par\vspace{-0.8em}%
}
\makeatother

% \blank[width] is an inline gap to write a short answer in.
\newcommand{\blank}[1][2.6cm]{\underline{\hspace{#1}}}

% A boxed space for showing working.
\newtcolorbox{workspace}[1][3.4cm]{%
  colback=white, colframe=rulegray, boxrule=0.5pt, arc=5pt,
  height=#1, left=6pt, right=6pt, top=4pt, bottom=4pt,
  before skip=9pt, after skip=11pt}

% An empty tick box.
\newcommand{\tickbox}{\fbox{\rule{0pt}{1.15em}\rule{1.15em}{0pt}}}

% A middot separator.
\newcommand{\sep}{\,\textperiodcentered\,}

% -- Section and question headings -------------------------------------------
% Light tint plus a fat coloured spine on the left. No solid dark banner: the
% colour reads clearly but barely costs the printer anything.
% needspace stops a heading stranding alone at the foot of a page. Kept modest:
% too large and it pushes headings over, leaving a worse hole than it prevents.
\newcommand{\hwsection}[2][cbteal]{%
  \needspace{8\baselineskip}%
  \par\vspace{13pt}%
  \noindent
  \begin{tcolorbox}[enhanced, nobeforeafter, width=\textwidth,
      colback=#1!7, colframe=#1!7, arc=5pt,
      borderline west={5pt}{0pt}{#1},
      left=14pt, right=10pt, top=7pt, bottom=7pt]
    {\large\bfseries\color{#1}#2}
  \end{tcolorbox}%
  \par\vspace{9pt}}

% \qhead[colour]{A2}{Moving along the line} -- a question heading.
\newcommand{\qhead}[3][cbteal]{%
  \needspace{6\baselineskip}%
  \par\vspace{11pt}%
  \noindent{\large\bfseries\color{#1}#2}\quad{\large\bfseries #3}\par\vspace{4pt}}

% -- Coloured boxes ----------------------------------------------------------
% Shared look: rounded, light fill, coloured rule, coloured title text on a
% pale title strip rather than a dark one.
\tcbset{hwbox/.style={
  enhanced, breakable, arc=6pt, boxrule=1pt,
  left=11pt, right=11pt, top=8pt, bottom=8pt,
  fonttitle=\bfseries, attach boxed title to top left={xshift=12pt, yshift=-8pt},
  boxed title style={arc=4pt, boxrule=1pt, left=7pt, right=7pt, top=2pt, bottom=2pt},
  top=13pt}}

% Orange: things to remember. Matches the deck's "Write This Down".
\newtcolorbox{remember}[1][Remember from the lesson]{hwbox,
  colback=cborange!6, colframe=cborange, coltitle=cborange,
  colbacktitle=white, title={#1}}

% Blue: English help. The barrier in this class is the wording, not the maths.
\newtcolorbox{wordhelp}[1][Word help]{hwbox,
  colback=cbblue!6, colframe=cbblue, coltitle=cbblue,
  colbacktitle=white, title={#1}}

% Purple: an activity or instruction.
\newtcolorbox{activity}[1][How to do this homework]{hwbox,
  colback=cbpurple!6, colframe=cbpurple, coltitle=cbpurple,
  colbacktitle=white, title={#1}}

% Green: your turn / a friendly nudge.
\newtcolorbox{yourturn}[1][Your turn]{hwbox,
  colback=cbgreen!7, colframe=cbgreen, coltitle=cbgreen!75!black,
  colbacktitle=white, title={#1}}

% -- Shorthands --------------------------------------------------------------
\newcommand{\dC}{\,^{\circ}\mathrm{C}}          % use inside maths: $-8\dC$
\newcommand{\key}[1]{{\color{cborange}\textbf{#1}}}
\newcommand{\eng}[1]{{\color{cbblue}\textbf{#1}}}

\setlist[enumerate]{itemsep=8pt, topsep=5pt, leftmargin=*}
\setlist[itemize]{itemsep=4pt, topsep=4pt, leftmargin=*}
\renewcommand{\arraystretch}{1.45}

% -- An expression the student is asked to work from -------------------------
\newcommand{\expr}[1]{%
  \tcbox[on line, colback=cbgreen!10, colframe=cbgreen!55, boxrule=1pt, arc=4pt,
         left=9pt, right=9pt, top=4pt, bottom=4pt]{\large\bfseries $#1$}}

% -- Running head and foot ---------------------------------------------------
% The packet overrides \fancyhead[L] with its own title.
\pagestyle{fancy}
\fancyhf{}
\fancyhead[L]{\small\color{black!55}Year 7\sep Homework}
\fancyhead[R]{\small\color{black!55}Name: \blank[3.4cm]}
\fancyfoot[C]{\small\color{black!55}Page \thepage}
\renewcommand{\headrulewidth}{0pt}
\renewcommand{\headrule}{\vspace{2pt}\hbox to\headwidth{\color{cbteal!45}\leaders\hrule height 1.2pt\hfill}}



% -- Per-packet running head -------------------------------------------------
\fancyhead[L]{\small\color{black!55}Year 7\sep Extra Homework}

\begin{document}

% ===========================================================================
% COVER BLOCK
% ===========================================================================
\thispagestyle{fancy}

\begin{tcolorbox}[enhanced, colback=cbgreen!8, colframe=cbgreen!8, arc=7pt,
                  borderline west={6pt}{0pt}{cbgreen},
                  left=18pt, right=14pt, top=14pt, bottom=14pt]
  {\color{cbgreen!75!black}\bfseries YEAR 7\sep UNIT 1\sep EXTRA HOMEWORK}\\[6pt]
  {\color{cbgreen!70!black}\Huge\bfseries The Two That Do the Work}\\[10pt]
  {\color{black!70}Multiplication and division\sep the engine room of Unit 1}
\end{tcolorbox}

\vspace{10pt}
\noindent
\begin{tabularx}{\textwidth}{@{}X X@{}}
  \textbf{Name:} \blank[5.6cm] & \textbf{Class:} \blank[5.6cm] \\
  \textbf{Date given:} \blank[4.6cm] & \textbf{Date due:} \blank[5.0cm] \\
\end{tabularx}

\vspace{8pt}
\begin{activity}[Why this sheet exists]
  Everything in Unit 1 sits on top of two operations. A \key{factor} is found by
  \key{dividing}, factor \key{pairs} by \key{multiplying}, and the \key{HCF}
  simplifies a fraction by \key{dividing} top and bottom. So the faster and more
  surely you can multiply and divide, the easier the whole unit gets. That is all
  this sheet is: \key{extra practice on the two operations that do the work.}

  \medskip
  \eng{Do it without a calculator}, and show your working --- working you can see
  is working you can check.
\end{activity}

% ===========================================================================
% E1 -- TABLES
% ===========================================================================
\hwsection[cbgreen]{Part 1 \textperiodcentered\ The facts worth knowing by heart}

\qhead[cbgreen]{E1}{Fill in the missing number}

\noindent These are the ones that come up again and again in factors and
divisibility. Aim to know them without counting.

\vspace{10pt}
\begin{multicols}{3}
\begin{enumerate}[label=(\alph*), itemsep=9pt]
  \item $6 \times 7 = \blank[1.3cm]$
  \item $7 \times 8 = \blank[1.3cm]$
  \item $8 \times 9 = \blank[1.3cm]$
  \item $9 \times 6 = \blank[1.3cm]$
  \item $7 \times 7 = \blank[1.3cm]$
  \item $8 \times 8 = \blank[1.3cm]$
  \item $12 \times 7 = \blank[1.3cm]$
  \item $11 \times 9 = \blank[1.3cm]$
  \item $9 \times 9 = \blank[1.3cm]$
  \item $12 \times 12 = \blank[1.3cm]$
  \item $8 \times 12 = \blank[1.3cm]$
  \item $9 \times 12 = \blank[1.3cm]$
  \item $7 \times 9 = \blank[1.3cm]$
  \item $11 \times 11 = \blank[1.3cm]$
  \item $7 \times 11 = \blank[1.3cm]$
\end{enumerate}
\end{multicols}

% ===========================================================================
% E2 -- MULTIPLICATION
% ===========================================================================
\qhead[cbgreen]{E2}{Multiply}

\noindent Set each one out properly and show your working in the box.

\vspace{6pt}
\noindent\textbf{(a)} Two digits by one digit.

\vspace{4pt}
\noindent
\begin{tabularx}{\textwidth}{|X|X|X|X|}
  \hline
  $34 \times 6$ & $47 \times 8$ & $58 \times 7$ & $93 \times 4$ \\
  \rule[-2.5em]{0pt}{3.1em} & & & \\
  \hline
\end{tabularx}

\vspace{12pt}
\noindent\textbf{(b)} Three digits by one digit.

\vspace{4pt}
\noindent
\begin{tabularx}{\textwidth}{|X|X|X|}
  \hline
  $236 \times 7$ & $418 \times 6$ & $364 \times 8$ \\
  \rule[-3.0em]{0pt}{3.8em} & & \\
  \hline
\end{tabularx}

\vspace{12pt}
\noindent\textbf{(c)} Two digits by two digits.

\vspace{4pt}
\noindent
\begin{tabularx}{\textwidth}{|X|X|X|}
  \hline
  $23 \times 14$ & $36 \times 25$ & $48 \times 17$ \\
  \rule[-3.6em]{0pt}{4.4em} & & \\
  \hline
\end{tabularx}

% ===========================================================================
% E3 -- DIVISION, EXACT
% ===========================================================================
\hwsection[cbgreen]{Part 2 \textperiodcentered\ Division}

\qhead[cbgreen]{E3}{Divide --- these all come out exactly}

\begin{remember}[A check that costs five seconds]
  When you finish a division, \key{multiply your answer back}. If
  $72 \div 6 = 12$, then $12 \times 6$ must come to 72. If it does not, you have
  just caught your own mistake.
\end{remember}

\vspace{4pt}
\begin{multicols}{2}
\begin{enumerate}[label=(\alph*), itemsep=13pt]
  \item $72 \div 6 = \blank[2.2cm]$
  \item $96 \div 4 = \blank[2.2cm]$
  \item $91 \div 7 = \blank[2.2cm]$
  \item $126 \div 9 = \blank[2.2cm]$
  \item $156 \div 12 = \blank[2.2cm]$
  \item $245 \div 5 = \blank[2.2cm]$
  \item $372 \div 4 = \blank[2.2cm]$
  \item $651 \div 3 = \blank[2.2cm]$
\end{enumerate}
\end{multicols}

% ===========================================================================
% E4 -- DIVISION WITH REMAINDERS
% ===========================================================================
\qhead[cbgreen]{E4}{Divide --- these do not}

\begin{wordhelp}[Word help --- the words that describe what is left]
  \begin{tabularx}{\linewidth}{@{}l@{\hspace{1.4em}}X@{}}
    \eng{exactly}          & nothing is left over. $72 \div 6 = 12$ exactly. \\
    \eng{remainder}        & what is left over when it does not come out exactly.
                             We write $53 \div 6 = 8$ r 5. \\
    \eng{goes into}        & ``6 goes into 24 four times'' means $24 \div 6 = 4$. \\
    \eng{left over}        & the everyday English for a remainder. \\
  \end{tabularx}
\end{wordhelp}

\noindent Write each answer as a whole number and a remainder, like this:
\quad $53 \div 6 = 8$ r 5.

\vspace{8pt}
\begin{multicols}{2}
\begin{enumerate}[label=(\alph*), itemsep=13pt]
  \item $74 \div 8 = \blank[2.8cm]$
  \item $100 \div 7 = \blank[2.8cm]$
  \item $85 \div 9 = \blank[2.8cm]$
  \item $137 \div 5 = \blank[2.8cm]$
  \item $191 \div 6 = \blank[2.8cm]$
  \item $250 \div 8 = \blank[2.8cm]$
\end{enumerate}
\end{multicols}

% ===========================================================================
% E5 -- FACT FAMILIES
% ===========================================================================
\hwsection[cbgreen]{Part 3 \textperiodcentered\ Putting them together}

\qhead[cbgreen]{E5}{Find the missing number}

\begin{remember}[Multiplying and dividing undo each other]
  If $7 \times 9 = 63$, then $63 \div 7 = 9$ \key{and} $63 \div 9 = 7$. Every
  multiplication fact is also \key{two} division facts --- which is why Part 1 is
  worth knowing by heart.
\end{remember}

\vspace{4pt}
\begin{multicols}{2}
\begin{enumerate}[label=(\alph*), itemsep=12pt]
  \item $7 \times \blank[1.4cm] = 63$
  \item $\blank[1.4cm] \times 8 = 96$
  \item $54 \div \blank[1.4cm] = 6$
  \item $\blank[1.4cm] \div 7 = 8$
  \item $6 \times \blank[1.4cm] = 84$
  \item $\blank[1.4cm] \times 9 = 117$
  \item $132 \div \blank[1.4cm] = 11$
  \item $\blank[1.4cm] \div 4 = 25$
\end{enumerate}
\end{multicols}

% ===========================================================================
% E6 -- USE IT
% ===========================================================================
\qhead[cbgreen]{E6}{Use it: is it a factor?}

\begin{remember}[Straight out of Homework 3]
  A number is a \key{factor} of another if it divides into it \key{exactly}.
  So the way to find out is to \key{do the division and look at the remainder}.
\end{remember}

\noindent Do each division, then circle \textbf{yes} or \textbf{no}.
\key{Two of these you have already worked out on this sheet} --- look back rather
than starting again.

\vspace{10pt}
\noindent
\begin{tabularx}{\textwidth}{|X|p{4.4cm}|p{3.4cm}|}
  \hline
  \textbf{Question} & \textbf{The division} & \textbf{Is it a factor?} \\
  \hline
  Is 6 a factor of 84?  \rule[-0.8em]{0pt}{2.3em} & & yes \quad / \quad no \\ \hline
  Is 8 a factor of 92?  \rule[-0.8em]{0pt}{2.3em} & & yes \quad / \quad no \\ \hline
  Is 9 a factor of 126? \rule[-0.8em]{0pt}{2.3em} & & yes \quad / \quad no \\ \hline
  Is 7 a factor of 100? \rule[-0.8em]{0pt}{2.3em} & & yes \quad / \quad no \\ \hline
\end{tabularx}

% ===========================================================================
% E7 -- WORD PROBLEMS
% ===========================================================================
\qhead[cbgreen]{E7}{Word problems}

\begin{wordhelp}[Word help --- which operation does the sentence want?]
  {\renewcommand{\arraystretch}{1.15}%
  \begin{tabularx}{\linewidth}{@{}l@{\hspace{1.4em}}X@{}}
    \eng{each} / \eng{every}        & usually multiply. 6 boxes with 24 in
                                      \textbf{each} is $6 \times 24$. \\
    \eng{share equally between}     & divide. \\
    \eng{how many are left}         & the remainder. \\
    \eng{altogether} / \eng{in total} & the whole amount, once you have finished. \\
  \end{tabularx}}
\end{wordhelp}

\noindent Write the calculation first, then the answer.

\begin{enumerate}[label=\textbf{\arabic*.}, itemsep=11pt]

  \item \textbf{The markers.} Mr Bowen buys \textbf{6 boxes} of markers. There are
        \textbf{24 markers} in each box. How many markers does he have altogether?
        \begin{workspace}[2.5cm]\end{workspace}

  \item \textbf{The paper.} Mr Bowen shares \textbf{96 sheets of paper} equally
        between \textbf{8 students}. How many sheets does each student get?
        \begin{workspace}[2.5cm]\end{workspace}

  \item \textbf{The pizzas.} Mr Bowen orders \textbf{8 identical pizzas} for the
        class. Each pizza is cut into \textbf{12 slices}. There are \textbf{29
        students}. Every student gets the same whole number of slices, and Mr Bowen
        eats the rest. How many slices does each student get, and how many does
        Mr Bowen eat?
        % The two-step one, and the last thing on the sheet: spend the leftover
        % space at the foot of the page on room to write rather than on nothing.
        \begin{workspace}[4.2cm]\end{workspace}

\end{enumerate}

% ===========================================================================
% ANSWER KEY (teacher copy only)
% ===========================================================================
\ifkey
\newpage
\hwsection[cbred]{Answer Key --- Teacher Copy}

\linespread{1.0}\selectfont
\setlist[enumerate]{itemsep=2pt, topsep=3pt, leftmargin=*}
\setlist[itemize]{itemsep=2pt, topsep=3pt, leftmargin=*}

\qhead[cbred]{E1}{Fill in the missing number}
\noindent (a) 42 \quad (b) 56 \quad (c) 72 \quad (d) 54 \quad (e) 49 \quad (f) 64
\quad (g) 84 \quad (h) 99 \\
(i) 81 \quad (j) 144 \quad (k) 96 \quad (l) 108 \quad (m) 63 \quad (n) 121
\quad (o) 77

\smallskip
\noindent{\small Time this one rather than marking it: if a student can do all
fifteen in under two minutes, the tables are not what is slowing them down in
Unit 1 and Part 2 is where to look instead. The three worth re-drilling if they
go wrong are (b) $7 \times 8$, (c) $8 \times 9$ and (m) $7 \times 9$ --- those
three account for most of the errors in factor lists.}

\qhead[cbred]{E2}{Multiply}
\noindent (a) $34 \times 6 = 204$ \quad $47 \times 8 = 376$ \quad
$58 \times 7 = 406$ \quad $93 \times 4 = 372$ \\
(b) $236 \times 7 = 1652$ \quad $418 \times 6 = 2508$ \quad $364 \times 8 = 2912$ \\
(c) $23 \times 14 = 322$ \quad $36 \times 25 = 900$ \quad $48 \times 17 = 816$

\smallskip
\noindent{\small The error to look for is carrying, not the tables: in (b),
$418 \times 6$ needs a carry at every column. In (c), check they have written
\emph{two} rows of working and added them --- a single row means they have
multiplied by the units digit only, and that shows up as an answer roughly a tenth
of the right size.}

\qhead[cbred]{E3}{Divide --- these all come out exactly}
\noindent (a) 12 \quad (b) 24 \quad (c) 13 \quad (d) 14 \quad (e) 13 \quad
(f) 49 \quad (g) 93 \quad (h) 217

\smallskip
\noindent{\small (c) $91 \div 7$ is the one that surprises people --- 91 looks
prime and is not. (h) $651 \div 3$ is the longest and is a good one to do at the
board. If a student's answers are right but slow, the multiply-back check in the
orange box is what to insist on next time, not more questions.}

\qhead[cbred]{E4}{Divide --- these do not}
\noindent (a) $74 \div 8 = 9$ r 2 \quad (b) $100 \div 7 = 14$ r 2 \quad
(c) $85 \div 9 = 9$ r 4 \\
(d) $137 \div 5 = 27$ r 2 \quad (e) $191 \div 6 = 31$ r 5 \quad
(f) $250 \div 8 = 31$ r 2

\smallskip
\noindent{\small Two errors to separate here. A remainder \emph{bigger than the
divisor} (writing $74 \div 8 = 8$ r 10) means the quotient is one too small and
they have stopped early --- that is a method error. A remainder that is simply
wrong by one or two is arithmetic. Only the first needs re-teaching. \\
Nobody should be writing a decimal here; if they do, it is worth saying that the
remainder is the useful form in this unit, because the remainder is what decides
whether something is a factor.}

\qhead[cbred]{E5}{Find the missing number}
\noindent (a) 9 \quad (b) 12 \quad (c) 9 \quad (d) 56 \quad (e) 14 \quad
(f) 13 \quad (g) 12 \quad (h) 100

\smallskip
\noindent{\small The four with the gap on the left of a division --- (d) and (h)
--- are the hard ones, because the missing number is the \emph{big} one and you
have to multiply to find it. Expect students to divide out of habit and get
$8 \div 7$. The fix is the orange box: say the fact family out loud. \\
(f) $13 \times 9 = 117$ is outside the tables on purpose; it is there to be worked
out rather than recalled.}

\qhead[cbred]{E6}{Use it: is it a factor?}
\noindent $84 \div 6 = 14$ --- \textbf{yes} \\
$92 \div 8 = 11$ r 4 --- \textbf{no} \\
$126 \div 9 = 14$ --- \textbf{yes} (already done in E3d) \\
$100 \div 7 = 14$ r 2 --- \textbf{no} (already done in E4b)

\smallskip
\noindent{\small The two repeats are deliberate. A student who does all four from
scratch has not read the instruction; a student who spots them has done the thing
this sheet is really for, which is noticing that the same division keeps turning
up. Ask who found them before you give the answers. \\
Note that rows 3 and 4 both give 14 --- one with nothing left over and one with a
remainder of 2. That pair is the whole idea of a factor in two lines, and it is
worth pointing at.}

\qhead[cbred]{E7}{Word problems}
\begin{enumerate}[label=\arabic*., itemsep=2pt]
  \item $6 \times 24 = \mathbf{144}$ markers.
  \item $96 \div 8 = \mathbf{12}$ sheets each.
  \item $8 \times 12 = 96$ slices altogether. $96 \div 29 = 3$ r 9, so each student
        gets \textbf{3 slices} and Mr Bowen eats \textbf{9}.
\end{enumerate}
\noindent{\small Problem 3 is the deadpan one. Read it completely straight and do
not explain the joke --- the maths is correct, which is the point. It is also the
only two-step question on the sheet: multiply first to find the total, then divide
with a remainder. A student who divides 8 by 29, or 12 by 29, has missed the first
step; send them back to the first sentence rather than to the arithmetic.}

\fi

\end{document}
